% ch2 C.1–C.2 (原书页 77–81 / PDF p092–p096)

% 衔接说明：本文件含大节 C 的开头与小节 C.1（引言）、C.2（弱束缚杂质态）。
% 文件最顶部（%--B4-TAIL-- 标记之下）为前一小节 B.4 的收尾文字（原书 p.77 顶部、
% 大节 C 标题之前的四段），装配时应移至 ch2_B4b.tex 末尾拼接。

% ---- 原书 p.77 (PDF p092)：至大节 C 标题之前为 B.4 收尾 ----

%--B4-TAIL--
%JOIN%
4s 将开始再度变空。）然而在 Cr 和 Mn，以及在某种程度上在 V 中，晶格开始变得远比 3d 电子的 $\kappa$ 所要求的稀疏，于是“3d 带”作为一个实体——有别于那些恰好在核附近才带有 3d 性质的单纯 O.P.W.——开始成形。

然而，由于未填满的 4s-4p 壳层的存在，这里的行为与 p 壳层系列中的行为颇为不同。4s-4p 壳层总保持着一定的占据，随着越来越多的 d 能级被占据、原子彼此相距相对更远，“O.P.W.”费米面便失去 3d 电子，回到只适合 4s-4p 电子的情形。于是在 V 与 Mn 之间的某处，d 带本身——作为由处于近似原子状态的电子构成的紧束缚带——得以成形；在此点之前，d 带与其他任何能带在定性上并无区别。

“内层”能带在这一点上的出现，以行为的巨大变化为标志，正如 O.P.W. 在 p 壳层的失效以分子结构与价结构的出现为标志一样。这里金属性得以保留，但内层电子的特征磁性开始显露——表现为 3d 行中实际的铁磁性与反铁磁性，以及 Pd 族和 Pt 族中的高磁化率与非典型行为。

在这一段题外话中，我试图说明：近满壳层中的空穴与近空壳层中的电子，为什么以及如何在其量子化学性质上表现得截然不同；也试图说明近自由电子图像在哪里、又为什么不能再指望奏效。我相信这样的分析以前从未有人做过。

\section{有微扰场存在时的单电子能带理论}

\subsection{引言}

在离开单电子理论这一主题之前，讨论当通常的周期场上叠加某种非周期微扰时固体中电子的行为，似乎是相当重要的。自然，对固体中处于完全平衡的自由电子所做的实验，无法告诉我们多少关于它们的信息——充其量，除总能量之外，我们至多能通过比热，或者（若运气好的话）通过 X 射线与光学光谱的测量，了解到态密度。为了研究能带结构，我们必须先微扰它。微扰可以由外加的电场和/或磁场构成——相对于原子尺度而言，它们实际上可视为恒定——也可以由周期势受到这种或那种扰动而产生的局域微扰势构成——杂质、位错（dislocation）、表面等等。

自然，真正微弱而局域的微扰——例如由与核磁矩的超精细相互作用、或由两种同位素核之间的差异所引起的微扰——可以用寻常的微扰论处理而无明显的复杂性。真正强的局域微扰——通常由结构中的间隙原子或缺失原子（空位）所引起——则无论在固体中还是在别处，都是一个本质上困难的问题。

% ---- 原书 p.78 (PDF p093) ----

另一方面，碰巧固体中人们感兴趣的最常见一类微扰，是这样一类：它们在总体上可以很强，却相对于原子尺度缓变。例如，均匀电场在哈密顿量中给出一个微扰项 $Eex$，即使 $E$ 很小，它在大晶体中也能变得要多大有多大：$x \to \infty$。均匀磁场在数值上同样可以很大：矢量势可取为 $\mathbf{r} \times \mathbf{H}/2$，它随位置矢量 $\mathbf{r}$ 线性增长。由位错、带电杂质、激子（exciton）、表面等等引起的势分布，也常常是缓变而幅度很大的。正是在这个极限下，“有效哈密顿量”（effective Hamiltonian）理论可以指望是有效的。

在这个理论中，所作的假定是：电子完全在单独一支能带的状态之间运动。这一假定的正当性我们稍后再讨论，但它显然足够合理：一个近乎恒定而平滑变化的势，可以使电子在能量几乎相同的状态之间绝热地来回移动，却不会引起向其他能带的真实跃迁。

我将首先在“有效哈密顿量”或“有效质量”（effective mass）近似下、并在上述假定之内，对电子的动力学作一简短讨论；然后再较深入地考察这一近似的某些背景。有不少文献对其中电子动力学有相当完整的讨论：Ziman (3) 论电场与磁场，Kohn 在《Seitzschrift》第 5 卷中论杂质态（impurity states）(44)，Wannier（文献 4 第 6 章）——我将多少遵循其进路——以及其他多种。Kittel 对“有效质量”的缘由有一段相当漂亮的定性讨论，我推荐作为配套读物（第 2 版，第 288 页），因为我这里不打算采用那种讲法。Wannier 的书在背景论证方面相当不错，不逊于任何参考文献，只有一部例外——它艰深然而完备：即《Seitzschrift》第 13 卷中的 E. I. Blount。

这个问题在如此长的时间里始终是一个活跃课题，本身就值得注意。相应的自由电子结果早在 1930 年就已知道——即 Landau 的抗磁性——其后不久，Jones 和 Zener 以及 Peierls 分别就电场情形和磁场情形，把它们推广到有效哈密顿量近似下的“准自由”电子 (46)。尽管如此，第一次试图体面地推导这些结果、并理解它们在何种意义上是近似的，来自 Wannier 1937 年关于 Wannier 函数的开创性论文。Seitz 书中的推导相当不完整。同样，这一近似的第一次失效由 Zener 于 1934 年——即 Zener 效应——以及 Houston 于 1940 年讨论过。然而，只有靠过去十年间 Luttinger、Adams、Wannier、Kohn、Blount 和 Gibson 的工作，人们才感到对这一问题达到了大概是最终的理解 (45)。有点令人惊讶的是，求解一个形式特别简单的单粒子 Schrödinger 方程竟然需要这一番功夫；而这些结果恰恰生动地表明，即便是自然界最简单的现象，其细节也可以何等复杂。

那么让我们假定：在缓变势存在时，电子确实保持在单独一支能带之内，其能量曲线为 $E_n(\mathbf{k})$——我们讨论的是第 $n$ 支带。

% ---- 原书 p.79 (PDF p094) ----

马上就引进 Wannier 函数（Wannier functions）$a_n(\mathbf{r} - \mathbf{R}_j)$ 的概念是很有教益的。这些函数由一个能带的 Bloch 函数
\begin{equation*}
b_n(\mathbf{k},\mathbf{r}) = \mathrm{e}^{i\mathbf{k}\cdot\mathbf{r}}\, u^n_{\mathbf{k}}(\mathbf{r})
\end{equation*}
出发，通过如下的 Fourier 分析得到：
\begin{equation*}
a_n(\mathbf{r} - \mathbf{R}_j) = N^{-1/2} \sum_{\mathbf{k}} \mathrm{e}^{-i\mathbf{k}\cdot\mathbf{R}_j}\, b_n(\mathbf{k},\mathbf{r}) \tag{19}
\end{equation*}
于是
\begin{equation*}
b_n = N^{-1/2} \sum_{\mathbf{R}_j} \mathrm{e}^{i\mathbf{k}\cdot\mathbf{R}_j}\, a_n(\mathbf{R}_j)
\end{equation*}

在紧束缚带中，$a_n$ 显然就是原来的紧束缚函数，因为此时 $u^n_{\mathbf{k}}(\mathbf{r})$ 与 $\mathbf{k}$ 无关，而且 $\sum_{\mathbf{k}} \to \delta(\mathbf{r} - \mathbf{R}_j)$；否则，$a_n$ 是一组适当正交且归一的函数，多少局域在格点 $\mathbf{R}_j$ 附近。在单带理论中，$a_n(\mathbf{r} - \mathbf{R}_j)$ 代表所能达到的最接近位置本征函数的东西；对自由电子，它们对应于 $\delta(\mathbf{r} - \mathbf{R})$。让我们定义一个算符 $\mathbf{R}$，使得 $\mathbf{R}\, a_n(\mathbf{r} - \mathbf{R}_j) = \mathbf{R}_j\, a_n(\mathbf{r} - \mathbf{R}_j)$。在 $\mathbf{k}$ 只占据第一区、而 $\mathbf{R}$ 以一组分立的格点位置为本征值这两个限制之下，它们代表一对彼此共轭的变量，类似于自由电子理论中的 $\mathbf{p}$ 和 $\mathbf{r}$，即
% 校注：按本页下文的推导（$i\,\partial/\partial\mathbf{R}$ 等价于 $\mathbf{k}$），
% 此对易子似应为 $[\mathbf{k},\mathbf{R}] = i$；原书印作 “= 1”，照录。
\begin{equation*}
[\mathbf{k}, \mathbf{R}] = 1
\end{equation*}
这是因为——按定义——
\begin{align*}
i\,\frac{\partial}{\partial\mathbf{R}}\; a_n(\mathbf{r} - \mathbf{R}_j) &= i\,\frac{\partial}{\partial\mathbf{R}_j}\; a_n(\mathbf{r} - \mathbf{R}_j) \\
&= \sum_{\mathbf{k}} \mathbf{k}\, b_n(\mathbf{k},\mathbf{r})\, \mathrm{e}^{i\mathbf{k}\cdot\mathbf{R}_j}
\end{align*}
［把 (19) 代入 $a_n$］。于是 $i(\partial/\partial\mathbf{R})$ 等价于 $\mathbf{k}$，反之亦然。现在，若 $V(\mathbf{r})$ 缓变，我们可以近似
\begin{equation*}
V(\mathbf{r})\, a_n(\mathbf{r} - \mathbf{R}_j) \simeq V(\mathbf{R}_j)\, a_n(\mathbf{r} - \mathbf{R}_j) = V(\mathbf{R})\, a_n(\mathbf{r} - \mathbf{R}_j)
\end{equation*}

% ---- 原书 p.80 (PDF p095) ----

一般说来，对于缓变的磁微扰或电微扰，可以有效地假定 $V(\mathbf{r})$ 实际上就等于 $V(\mathbf{R})$。于是，一个近似哈密顿量为
\begin{align*}
\mathcal{H} &= E_n(\mathbf{k}) + V(\mathbf{r}) \simeq \mathcal{H}_{\text{eff}} = E_n(\mathbf{k}) + V(\mathbf{R}) \\
\mathcal{H}_{\text{eff}} &= E_n\!\left(i\,\frac{\partial}{\partial\mathbf{R}}\right) + V(\mathbf{R}) \quad\text{或}\quad = E_n(\mathbf{k}) + V\!\left(-i\,\frac{\partial}{\partial\mathbf{k}}\right)
\end{align*}
在同样的限制下，显然可以同样不困难地证明：当存在使
\begin{equation*}
E_{\text{kin}} \to -\hbar^2\left(\nabla - \frac{e\mathbf{A}}{i\hbar c}\right)^{\!2}\Big/2m
\end{equation*}
的磁微扰时，把 $\mathbf{p}$ 换成 $\hbar\mathbf{k}$、把 $A(\mathbf{r})$ 换成 $A(\mathbf{R})$ 是正确的：
\begin{equation*}
\mathcal{H}_{\text{eff}} \simeq E_n\!\left(\mathbf{k} - \frac{e\mathbf{A}(\mathbf{R})}{\hbar c}\right) + V(\mathbf{R})
\end{equation*}

\subsection{弱束缚杂质态}

这一形式理论所能应用的三类问题中，最简单的是半导体中带电杂质原子的问题，例如 Si 中的 P (44)。在 Si 中，P 以替位方式占据一个 Si 格位，但当然带有一个额外的正电荷；它是一个“施主”（donor）。其结果是：在足够低的温度下，一个额外的价电子被束缚在过剩正电荷的区域中。

碰巧，这样得到的势相当弱，相应的波函数相当长程，这既因为 Si 的介电常数（dielectric constant）$\kappa$ 相当大（势为 $V = 2e^2/\kappa r$），也因为有效质量很小。因此可以采用“有效质量”或“有效哈密顿量”近似。就 Si 中的价电子而言，能量面 $E_n(\mathbf{k})$ 由一组六个“谷”（valleys）组成，这些谷沿六个 100 方向位于点 $(k_o, 0, 0)$ 处，其中 $k_o \cong 2\pi/a$（0.85）。在这些谷附近
\begin{equation*}
E_n(\mathbf{k}) \simeq \frac{\hbar^2 (k_x - k_o)^2}{2m^*_{\parallel}} + \frac{\hbar^2 (k_y^2 + k_z^2)}{2m^*_{\perp}}
\end{equation*}
于是，让我们把其中一个谷中的试探波函数（trial wave function）写成
\begin{equation*}
f(\mathbf{r}) = F(\mathbf{R})\, b_{k_o} = F(\mathbf{R}) \sum_{\mathbf{R}_j} \mathrm{e}^{i k_o \cdot \mathbf{R}_j}\, a_n(\mathbf{r} - \mathbf{R}_j) \tag{20}
\end{equation*}
并注意到
\begin{equation*}
k_x\, b_{k_o} = i\,\frac{\partial}{\partial R_x}\, b_{k_o} = k_o\, b_{k_o} \qquad\qquad k_y\, b_{k_o} = k_z\, b_{k_o} = 0
\end{equation*}
于是
\begin{align*}
(k_x - k_o)^2 F(\mathbf{R})\, b_{k_o} ={}& -\left[\frac{\partial^2}{\partial R_x^2}\, F(\mathbf{R})\right] b_{k_o} \\
&- 2i\left[\frac{\partial}{\partial R_x}\, F(\mathbf{R})\right] k_o\, b_{k_o} + F(\mathbf{R})\, k_o^2\, b_{k_o} \\
&- 2k_o \left\{ -\left[i\,\frac{\partial}{\partial R_x}\, F(\mathbf{R})\right] b_{k_o} + F(\mathbf{R})\, k_o\, b_{k_o} \right\} \\
&+ k_o^2 F(\mathbf{R})\, b_{k_o} = -\left[\frac{\partial^2}{\partial R_x^2}\, F(\mathbf{R})\right] b_{k_o}
\end{align*}
这样，$b_{k_o}$ 便可以从整个波动方程中作为公因子提出，我们得到的简而言之就是一个各向异性的氢原子波动方程：

% ---- 原书 p.81 (PDF p096) ----

\begin{equation*}
-\frac{\hbar^2}{2m^*_{\parallel}}\, \frac{\partial^2 F(\mathbf{R})}{\partial R_x^2} - \frac{\hbar^2}{2m^*_{\perp}} \left(\frac{\partial^2}{\partial R_y^2} + \frac{\partial^2}{\partial R_z^2}\right) F(\mathbf{R}) - \frac{2e^2}{\kappa R}\, F(\mathbf{R}) = E\, F(\mathbf{R}) \tag{21}
\end{equation*}

这样得到的波函数，其半径大致为 $\hbar^2 \kappa / \langle m^* \rangle_{\text{av}}\, e^2$，在 Si 中约为 50 个 Bohr 半径；其中 $\kappa$ 为 12，而 $\langle m^* \rangle_{\text{av}}$［所应使用的适当平均，必须由式 (21) 的数值解得出］约为 $\frac{1}{4} m_e$。

在这个近似下，有六个简并的波函数，能量面上的每个谷给出一个。在中心原子处，有效质量理论绝不可能真正有效，理论的这一失效造成一个微扰，把这一简并分裂开；最低的波函数是全部六个函数的对称化线性组合，因而在中心原子处具有相当大的振幅，其能量也离式 (21) 的解相当远。Kohn 和 Luttinger 曾极其详细地考察过对有效质量理论的偏离。在对此效应作出修正之后，Kohn 和 Schechter 算出了 $F(\mathbf{R})$ 相当剧烈的变化，它来源于式 (20) 的六个函数之间的干涉。Feher 的经典实验测量了各种不同邻近格点处的振幅，对这一理论给出了定性上良好、但定量上尚不完善的支持 (47)。关于这些有效质量波函数的进一步细节，在这里就不适宜多谈了。
